\section*{Bab \ref{ch:g11:reaction-energy} --- Energi Reaksi: Pembakaran dan Energi Ikatan}

\begin{solution}{exo:g11:reaction-energy:1}
Kayu yang \omterm{def:g4:what-a-fire-needs:burning}{terbakar}: \omterm{def:g11:reaction-energy:exothermic}{eksoterm}. Es yang mencair: \omterm{def:g11:reaction-energy:exothermic}{endoterm} (es menyerap
panas dari lingkungannya, walaupun pencairan bukan \omterm{def:g7:chemical-reactions:reaction}{reaksi kimia}).
Kantong dingin instan: \omterm{def:g11:reaction-energy:exothermic}{endoterm}. Penghangat tangan: \omterm{def:g11:reaction-energy:exothermic}{eksoterm}. Pembuatan
gula di bawah sinar matahari: \omterm{def:g11:reaction-energy:exothermic}{endoterm} (energinya berasal dari cahaya).
\end{solution}

\begin{solution}{exo:g11:reaction-energy:2}
$Q = 3.0 \times 890.7 = \qty{2.7e3}{kJ}$, sekitar \qty{2.7}{MJ}.
\end{solution}

\begin{solution}{exo:g11:reaction-energy:3}
\omterm{def:g7:chemical-reactions:reactant}{Pereaksi}. Selisihnya dilepaskan ke lingkungan, sebagian besar sebagai
panas.
\end{solution}

\begin{solution}{exo:g11:reaction-energy:4}
$n = 100 / 46.0 = \qty{2.17}{mol}$; $Q = 2.17 \times 1367.6 =
\qty{2.97e3}{kJ}$, sekitar \qty{3.0}{MJ}.
\end{solution}

\begin{solution}{exo:g11:reaction-energy:5}
Energi yang diperlukan untuk memutus satu \omterm{def:g10:the-mole:amount}{mol} ikatan semacam itu, dengan
\omterm{def:g7:atoms-and-molecules:molecule}{molekul} dan \omterm{def:g7:atoms-and-molecules:atom}{atomnya} berwujud gas. Pasangan \omterm{def:g9:inside-the-atom:nucleus}{elektron} bersama itulah yang
mengikat kedua \omterm{def:g7:atoms-and-molecules:atom}{atom}: menarik keduanya terpisah berarti melawan gaya
tarik itu, dan hal itu memerlukan energi.
\end{solution}

\begin{solution}{exo:g11:reaction-energy:6}
Diputus: $2 \times 436 + 498 = \qty{1370}{kJ}$; dibentuk:
$4 \times 464 = \qty{1856}{kJ}$; $E_r \approx 1370 - 1856 =
\qty{-486}{kJ}$: \omterm{def:g11:reaction-energy:exothermic}{eksoterm}.
\end{solution}

\begin{solution}{exo:g11:reaction-energy:7}
Etanol \ce{CH3-CH2-O-H} memiliki 5 \ce{C-H}, 1 \ce{C-C}, 1 \ce{C-O},
dan 1 \ce{O-H}. Diputus: $5 \times 415 + 345 + 350 + 464 + 3 \times 498 =
\qty{4728}{kJ}$. Dibentuk: $4 \times 741 + 6 \times 464 = \qty{5748}{kJ}$.
$E_r \approx \qty{-1020}{kJ}$, besarnya sekitar seperempat lebih kecil
daripada nilai terukur \qty{1367.6}{kJ}.
\end{solution}

\begin{solution}{exo:g11:reaction-energy:8}
Propana: $2219.2 / 0.0440 = \qty{5.04e4}{kJ/kg} = \qty{50.4}{MJ/kg}$.
Metana: $890.7 / 0.0160 = \qty{55.7}{MJ/kg}$. Keduanya sesuai dengan
diagram.
\end{solution}

\begin{solution}{exo:g11:reaction-energy:9}
$Q = 1500 \times 4.18 \times 85 = \qty{5.33e5}{J} = \qty{533}{kJ}$;
$n = 533 / 890.7 = \qty{0.598}{mol}$, yaitu $0.598 \times 16.0 =
\qty{9.6}{g}$ metana; jika setengah energinya hilang, \qty{19}{g}.
\end{solution}

\begin{solution}{exo:g11:reaction-energy:10}
Hidrogen > metana > propana > oktana > etanol.
$48.0 / 142.9 = \qty{0.34}{kg}$ hidrogen.
\end{solution}

\begin{solution}{exo:g11:reaction-energy:11}
Diputus: $946 + 3 \times 436 = \qty{2254}{kJ}$; dibentuk:
$6 \times 390 = \qty{2340}{kJ}$; $E_r \approx \qty{-86}{kJ}$: sedikit
\omterm{def:g11:reaction-energy:exothermic}{eksoterm}.
\end{solution}

\begin{solution}{exo:g11:reaction-energy:12}
Air: $250 \times 4.18 \times 20 = \qty{2.09e4}{J} = \qty{20.9}{kJ}$.
Etanol: $1.20 / 46.0 = \qty{0.0261}{mol}$, yang dapat melepaskan
$0.0261 \times 1367.6 = \qty{35.7}{kJ}$. Efisiensi
$20.9 / 35.7 \approx \qty{59}{\percent}$. Kehilangan energi: panas yang
terbawa \omterm{def:g4:air-a-mixture-of-gases:air}{udara}, panas yang menghangatkan kaleng dan penyangganya,
\omterm{def:g8:combustion-and-fuels:complete}{pembakaran tidak sempurna} (jelaga), sebagian etanol \omterm{def:g3:separating-mixtures:evaporate}{menguap} tanpa
\omterm{def:g4:what-a-fire-needs:burning}{terbakar}.
\end{solution}

\begin{solution}{exo:g11:reaction-energy:13}
\ce{2C8H18 + 25O2 -> 16CO2 + 18H2O}: 8 \omterm{def:g10:the-mole:amount}{mol} karbon dioksida,
\qty{352}{g}, per \qty{5470}{kJ}, jadi $352 / 5.470 = \qty{64.4}{g}$ per
megajoule. \ce{C3H8 + 5O2 -> 3CO2 + 4H2O}: \qty{132}{g} per
\qty{2219.2}{kJ}, jadi \qty{59.5}{g} per megajoule. Metana, dengan
\qty{49.4}{g}, melepaskan paling sedikit.
\end{solution}

\begin{solution}{exo:g11:reaction-energy:14}
Tabel memberikan \omterm{def:g11:reaction-energy:bond-energy}{energi ikatan} rata-rata, sedangkan ikatan \ce{C=O}
karbon dioksida lebih kuat daripada rata-rata; nilai terukurnya berlaku
untuk air cair, yang pengembunannya melepaskan energi lebih banyak
daripada reaksi gasnya; dan metode itu memperlakukan setiap ikatan
seolah-olah tidak bergantung pada ikatan di sekitarnya. Perkiraan itu
tetap berguna: perkiraan itu memberikan tanda yang benar (\omterm{def:g11:reaction-energy:exothermic}{eksoterm}) dan
orde besar yang benar tanpa pengukuran apa pun.
\end{solution}

\begin{solution}{exo:g11:reaction-energy:15}
Diputus: $4 \times 464 = \qty{1856}{kJ}$; dibentuk: $2 \times 436 + 498 =
\qty{1370}{kJ}$; $E_r \approx \qty{+486}{kJ}$: \omterm{def:g11:reaction-energy:exothermic}{endoterm}, jadi energi
harus diberikan. Membakar hidrogen itu kemudian mengembalikan energi
ini: hidrogen menyimpan energi yang dihasilkan di tempat lain, hidrogen
tidak menciptakannya.
\end{solution}

\begin{solution}{pb:g11:reaction-energy:1}
\textbf{1.} \ce{CH4 + 2O2 -> CO2 + 2H2O}.

\textbf{2.} \ce{C2H6O + 3O2 -> 2CO2 + 3H2O}.

\textbf{3.} \ce{2C8H18 + 25O2 -> 16CO2 + 18H2O}.

\textbf{4.} \ce{2H2 + O2 -> 2H2O}.

\textbf{5.} Hidrogen.

\textbf{6.} \qty{16.0}{g/mol}, \qty{46.0}{g/mol}, \qty{114.0}{g/mol},
\qty{2.0}{g/mol}.

\textbf{7.} Dalam \si{kJ} per \si{kg}, lalu \si{MJ/kg}: metana
$890.7 / 0.0160$, \qty{55.7}{MJ/kg}; etanol $1367.6 / 0.0460$,
\qty{29.7}{MJ/kg}; oktana $5470 / 0.1140$, \qty{48.0}{MJ/kg}; hidrogen
$285.8 / 0.0020$, \qty{142.9}{MJ/kg}.

\textbf{8.} Hidrogen > metana > oktana > etanol.

\textbf{9.} Hidrogen adalah gas yang sangat ringan: satu kilogram
hidrogen mengisi volume yang sangat besar pada tekanan normal, sehingga
hidrogen harus dimampatkan ke dalam tangki bertekanan tinggi yang berat.

\textbf{10.} Metana: 4 \ce{C-H}; oksigen: 2 \ce{O=O} diputus. Karbon
dioksida: 2 \ce{C=O}; air: $2 \times 2 = 4$ \ce{O-H} dibentuk.

\textbf{11.} $4 \times 415 + 2 \times 498 = \qty{2656}{kJ}$.

\textbf{12.} $2 \times 741 + 4 \times 464 = \qty{3338}{kJ}$.

\textbf{13.} $E_r \approx 2656 - 3338 = \qty{-682}{kJ}$, dibandingkan
\qty{-890.7}{kJ} terukur: besarnya sekitar \qty{23}{\percent} terlalu
kecil.

\textbf{14.} \omterm{def:g11:reaction-energy:bond-energy}{Energi ikatan} rata-rata (\ce{C=O} karbon dioksida lebih
kuat daripada rata-rata); air cair dalam pengukuran, gas dalam
perkiraan.

\textbf{15.} $1000 / 890.7 = \qty{1.123}{mol}$.

\textbf{16.} $1.123 \times 44.0 = \qty{49.4}{g}$.

\textbf{17.} Etanol: $1000 / 1367.6 = \qty{0.731}{mol}$, yang
menghasilkan \qty{1.462}{mol} karbon dioksida, \qty{64.3}{g}. Oktana:
$1000 / 5470 = \qty{0.183}{mol}$, yang menghasilkan \qty{1.463}{mol},
\qty{64.4}{g}.

\textbf{18.} Metana: metana memiliki \omterm{def:g7:atoms-and-molecules:atom}{atom} hidrogen paling banyak per
\omterm{def:g7:atoms-and-molecules:atom}{atom} karbon (4, dibandingkan 2.25 untuk oktana), dan membakar hidrogen
memberikan energi tanpa karbon dioksida.

\textbf{19.} Pada cara hidrogen itu dibuat: jika dibuat dari air dengan
listrik dari sumber terbarukan, hidrogen hampir tidak melepaskan karbon
dioksida; jika dibuat dari metana, karbon dioksida dilepaskan di pabrik,
bukan di knalpot.

\textbf{20.} Metana melepaskan sekitar \qty{49}{g} karbon dioksida per
megajoule.
\end{solution}
